\chapter{Toolchain}
\label{chap:toolchain}

\begin{table}[H]
\caption{History / Revision – Chapter 5}
\label{tab:history-ch1}
\centering
\begin{tabular}{llll}
\toprule
Version & Date & Author & Remarks \\
\midrule
1.0 & January 22, 2026 & Y. Sangale & Initial template draft \\
\bottomrule
\end{tabular}
\end{table}

\section{Introduction of the Toolchain}


The development of a modern System-on-Chip (SoC) requires a robust, automated, and developer-friendly toolchain that seamlessly integrates hardware design, firmware development, simulation, testing, and documentation.  
For the \textbf{RV64I\_RWU} project, a comprehensive toolchain was designed and extended to support the full development lifecycle of a 64-bit RISC-V processor core.

The toolchain is centered around a \textbf{modular CMake-based build system} and is complemented by a custom-developed \textbf{Visual Studio Code (VS Code) extension}.  
This combination enables both command-line–driven workflows for automation and graphical workflows for developer productivity. The primary goals of the toolchain are:

\begin{itemize}
    \item Reproducible and cross-platform builds
    \item Automated firmware compilation and hardware simulation
    \item Scalable regression testing using CTest
    \item Reduced onboarding effort for new developers
    \item Tight integration between software, hardware, and verification flows
\end{itemize}

By abstracting complex build and simulation commands behind well-defined interfaces, the toolchain significantly improves usability while maintaining full transparency and flexibility for advanced users.

\section{Technologies Used}

The toolchain is built using industry-standard tools and technologies to ensure portability, maintainability, and long-term scalability.

\subsection{Build and Automation Technologies}

\begin{itemize}
    \item \textbf{CMake ($\geq$ 3.20)}:  
    Used as the primary build system generator. A modular CMake architecture enables auto-discovery of RTL, firmware, and test files while maintaining a clean separation of concerns.
    
    \item \textbf{CTest}:  
    Integrated with CMake to provide automated unit, integration, firmware, and IP-level regression testing with support for parallel execution.
    
    \item \textbf{RISC-V GNU Toolchain}:  
    The \texttt{riscv64-unknown-elf-gcc} compiler is used for compiling firmware written in assembly and C for the RV64I architecture.
\end{itemize}

\subsection{Simulation and Verification Tools}

\begin{itemize}
    \item \textbf{Xilinx Vivado XSIM}:  
    Used for RTL simulation, supporting both console-based execution for automation and GUI-based execution with waveform visualization.
    
    \item \textbf{TCL Scripts}:  
    Employed to standardize waveform setup, simulation control, and batch execution in XSIM.
\end{itemize}

\subsection{Development and Integration Tools}

\begin{itemize}
    \item \textbf{Visual Studio Code}:  
    Serves as the primary development environment.
    
    \item \textbf{Node.js}:  
    Used to develop a custom VS Code extension that provides an integrated user interface for the toolchain.
    
    \item \textbf{Git}:  
    Used for version control and collaborative development.
    
    \item \textbf{Sphinx}:  
    Used for generating HTML and PDF documentation directly from the project source.
\end{itemize}

\section{Usage of the Toolchain}

The RV64I toolchain supports two primary usage modes: a terminal-based workflow and an IDE-based workflow using the VS Code extension.

\subsection{Terminal-Based Usage}

The terminal workflow provides full control over the build and simulation process and is well suited for automation, scripting, and continuous integration environments.

\subsubsection{Configuration and Build}

\begin{verbatim}
cmake -S . -B build -G "MinGW Makefiles"
cmake --build build
\end{verbatim}

This step configures the project, detects available tools, and builds all auto-discovered targets.

\subsubsection{Firmware Compilation}

\begin{verbatim}
cmake --build build --target fw_asm_instr06addi
\end{verbatim}

Firmware targets are generated automatically based on the source files present in the firmware directory.

\subsubsection{Simulation Execution}

\textbf{Console Mode:}
\begin{verbatim}
xsim instr06addi_snapshot -runall
\end{verbatim}

\textbf{GUI Mode with Waveforms:}
\begin{verbatim}
xsim instr06addi_snapshot --gui --tclbatch sim/config/xsim_cfg2.tcl
\end{verbatim}

\subsubsection{Regression Testing}

\begin{verbatim}
ctest --test-dir build -j
\end{verbatim}

CTest enables parallel execution of unit, integration, firmware, and IP-level tests.

\subsection{VS Code Extension Usage}

To improve usability and reduce the cognitive load of complex command sequences, a custom VS Code extension named \textbf{rwu-rv64i-toolchain} was developed.

\subsubsection{Setup Wizard}

The extension provides an interactive setup wizard that:
\begin{itemize}
    \item Detects installed tools (CMake, RISC-V GCC, XSIM, Git)
    \item Reports missing components
    \item Allows manual configuration of executable paths
\end{itemize}

This significantly simplifies first-time setup, especially on Windows systems.

\subsubsection{Build and Simulation Commands}

The extension exposes high-level commands via the Command Palette, including:
\begin{itemize}
    \item Configure and Build All
    \item Build Specific Target (compile only)
    \item Run Simulation (console mode)
    \item Run Simulation (GUI with waveforms)
\end{itemize}

Target names are automatically inferred from the currently opened file, reducing manual input errors.

\subsubsection{Regression Testing Support}

The extension integrates CTest commands, enabling:
\begin{itemize}
    \item Execution of specific tests
    \item Full library regression testing
    \item Parallel test execution
\end{itemize}

All test results are displayed directly in the VS Code terminal for immediate feedback.

\section{Further Development of the Toolchain}

While the current toolchain is fully functional, several extensions and improvements are planned.

\subsection{CMake Enhancements}

Future improvements to the CMake-based build system include:
\begin{itemize}
    \item Support for multiple simulators beyond XSIM
    \item Improved dependency visualization and reporting
    \item Enhanced cross-compilation support for additional RISC-V variants
    \item Integration with continuous integration (CI) pipelines
\end{itemize}

\subsection{VS Code Extension Enhancements}

Planned enhancements to the VS Code extension include:
\begin{itemize}
    \item Graphical dashboards for build and test status
    \item Integrated waveform viewers
    \item Per-target configuration profiles
    \item Deeper integration with documentation and architecture diagrams
\end{itemize}

These improvements aim to further reduce development time and improve the overall developer experience.
